% The construction rule in three pictures: a region, the one full wall with its
% single door, and the same rule applied again inside each half. Nothing is
% drawn cell by cell here -- the point is the rule, not an instance.
\begin{figure}[htbp]
  \centering
  \begin{subfigure}[b]{0.31\textwidth}\centering
    \begin{tikzpicture}[scale=0.42]
      \draw[region] (0,0) rectangle (8,5);
      \node[font=\scriptsize] at (4,2.5) {one region};
    \end{tikzpicture}
    \caption{A rectangular region, every cell reachable from every other.}
    \label{fig:rule-region}
  \end{subfigure}\hfill
  \begin{subfigure}[b]{0.31\textwidth}\centering
    \begin{tikzpicture}[scale=0.42]
      \draw[region] (0,0) rectangle (8,5);
      \draw[wall segment] (3,0) -- (3,2);
      \draw[wall segment] (3,3) -- (3,5);
      \draw[door] (3,2) -- (3,3);
      \node[doorgreen, font=\scriptsize, right] at (3.2,2.5) {door};
    \end{tikzpicture}
    \caption{One full wall across the longest side, pierced by exactly one
    door.}
    \label{fig:rule-one-cut}
  \end{subfigure}\hfill
  \begin{subfigure}[b]{0.31\textwidth}\centering
    \begin{tikzpicture}[scale=0.42]
      \draw[region] (0,0) rectangle (8,5);
      \draw[wall segment] (3,0) -- (3,2);
      \draw[wall segment] (3,3) -- (3,5);
      \draw[door] (3,2) -- (3,3);
      \draw[wall segment] (0,3) -- (1,3);
      \draw[wall segment] (2,3) -- (3,3);
      \draw[door] (1,3) -- (2,3);
      \draw[wall segment] (3,2) -- (6,2);
      \draw[wall segment] (7,2) -- (8,2);
      \draw[door] (6,2) -- (7,2);
    \end{tikzpicture}
    \caption{The same rule applied inside each of the two halves.}
    \label{fig:rule-recursion}
  \end{subfigure}
  \caption{The construction rule. Each application turns one region into two
  regions of the same kind joined by a single door, which is what makes this a
  divide-and-conquer problem and the decomposition a binary tree.}
  \label{fig:construction-rule}
\end{figure}
